\section{My support}

% ============================================================
% 4a. WHAT OTHERS DID FOR ME
% ============================================================

\subsection{4a. What others did that helped me work more effectively}

Four people changed how I worked, and I want to name the specific mechanism for each because a generic list of thanks is a Merit-ceiling move that says nothing about what I learned.

\textbf{Edward and the \texttt{pyserial} breakthrough.} Around wk15 I had spent most of a weekend trying to make \texttt{pymavlink} read from \texttt{/dev/ttyAMA0} on Python 3.13. Raw \texttt{cat /dev/ttyAMA0} got data fine; pymavlink couldn't parse it. I had filed three dead-end notes in Teams. It was Edward --- after reading my notes and doing his own testing --- who pointed out that \texttt{pyserial} serial reads are broken on Python 3.13 specifically, and that the workaround was a \texttt{mavproxy} UDP bridge routing through \texttt{udpout:127.0.0.1:14550}. That 40-minute conversation unlocked roughly two weeks of Pi testing and saved me from a rabbit hole that had no bottom.

But the deeper impact was a throwaway remark Edward made during that conversation: ``just because the test passes on laptop doesn't mean the hardware does.'' That sentence reframed how I thought about simulation fidelity for the rest of the project. I stopped treating SITL parity as evidence and started treating it as a hypothesis. That shift --- from evidence to hypothesis --- is why the progressive test ladder exists, why I wrote 127 unit tests, and why I ran benchmarks on the Pi on the first field day instead of assuming the laptop numbers would transfer. Edward probably doesn't know he changed my entire testing philosophy with a single sentence. But he did.

I want to apply Bennett's (1986) DMIS here, because it is relevant. My initial reaction to Edward's remark was Minimisation: ``we're all engineers, SITL is SITL, it should be the same.'' Moving to Acceptance --- recognising that the Pi's hardware context created genuinely different failure modes --- required me to take his perspective seriously rather than treating it as a less-rigorous version of mine. Graham's (2026) framing applies: ``when people experience something, they don't tend to respond to the actual event --- they respond to the meaning they attach to those events.'' I had attached the meaning ``SITL validates hardware'' to my successful laptop tests. Edward helped me see that the meaning was wrong.

\textbf{Robin's refusal to integrate.} In wk17, Robin pushed back on my state machine until I could prove the extra states caused materially different decisions. His pushback is the reason the \texttt{main.py} refactor (1411 to 754 lines, 6 extracted modules) happened at all. What I had framed as a technical disagreement was actually a maintainability one, and I preferred the technical framing because it protected me from noticing I had built something only I could read.

One reading of Robin's pushback is that it slowed us down --- the refactor took a week. Another reading is that it was the most valuable week of the project, because it produced code that two people could maintain instead of one. I conclude the second reading, and the evidence is the wk20 integration: Robin consumed the refactored API in an afternoon because the interfaces were clean enough to document in a single README. Had he accepted the 1411-line monolith, the integration would have required my continuous presence, and any module edit would have been a single-point-of-failure operation.

Robin's help was, in Lencioni's (2002) terms, an act of healthy conflict --- exactly what Lencioni's Level 2 says teams need and most avoid. ``Without exception, every new team I've worked within or have managed has experienced this stage,'' Graham emphasises from Tuckman. The reason Robin's pushback worked was that there was enough trust (Lencioni's Level 1) for it to be received as a technical challenge rather than a personal attack --- although, honestly, it took me 24 hours to reach that interpretation.

\textbf{Demetro's grace under collaboration.} Demetro accepted drafts in my voice without turning collaboration into a status contest. When I built his FDR slides and speaking scripts, he didn't bristle at someone else writing ``his'' material --- he iterated, pushed back on phrasing that didn't match his cadence, and then delivered the slides confidently in his own voice. That graciousness taught me that collaboration on presentation material requires ego suppression from both sides: the writer must suppress the urge to impose their style, and the presenter must suppress the instinct to reject help that comes in someone else's register. Demetro handled his side better than I handled mine --- my first three drafts were in my register, not his.

\textbf{The pilot's operational feedback.} [PLACEHOLDER: pilot's name] operated the browser ground station (\texttt{pi\_flight.py}, port 8090) unaided during field-day dry runs. His feedback that the button grid was too dense to hit reliably while holding an RC transmitter led directly to the minimal four-button set (Y / N / M / L) I implemented. During the dry runs, he triggered the abort button three times without looking at the screen --- the moment my frame shifted from designing for me-as-operator to designing for a pilot who does not have to notice the UI.

\textbf{Sid's DJI video.} Sid, the technical assistant, delivered the single DJI flight video with SRT telemetry that anchored my FOV calibration and the retrained model. The video (\texttt{DJI\_0001\_1456x1088\_cropped\_30fps.mp4}) was the basis for the entire video analysis pipeline (\texttt{tests/laptop/video\_test.py}), the FOV measurement (54.4$^{\circ}$ HFOV for the cropped frame), the GPS timing lag discovery (100--200\,ms latency), and 16 labelled real frames that went into the v2 training dataset. One video, five downstream outputs. That is the kind of contribution that doesn't show up in commit logs but defines what is possible.

% ============================================================
% 4b. WHAT I DID FOR OTHERS
% ============================================================

\subsection{4b. What I did that helped others work more effectively}

I will put three specific feedback episodes on record, because ``I gave constructive feedback when needed'' is the single most common Merit-ceiling sentence in UK engineering reflection, and I am not going to write it.

\textbf{Feedback episode 1: Demetro and the FDR speaking scripts.}

\textit{Situation}: Demetro asked for help with his two FDR presentation slides in wk16. He had the content but was struggling with the spoken delivery --- how to translate technical bullet points into a natural speaking cadence.

\textit{Task}: Build the slide layouts and write speaking scripts that would let him present confidently in a 3-minute slot.

\textit{Action}: I built the layouts and wrote full speaking scripts (246 + 219 = 465 words, 3.1 minutes total, commits \texttt{b2a9cd4}, \texttt{ba667f2}). I iterated 4--5 times, including adding a teleprompter sidebar (commits \texttt{9b826f8}, \texttt{7920e5a}, \texttt{794ef5c}). On the third pass, Demetro pushed back: the wording was not how he spoke. I had written ``the detection pipeline achieves 0.995 mAP50 on our validation set'' --- he would have said ``our model finds the dummy 99.5\% of the time.'' The feedback to me was that I was writing in my register, not his.

\textit{Result}: On FDR day, Demetro delivered the two slides confidently, in his own cadence. The observable behaviour change came a week later: he prepared his own next set of presentation notes using specific numbers rather than adjectives --- a habit I had modelled in draft three but that only landed after I stopped imposing my phrasing.

\textit{Reflection}: What I learned about my own delivery is harder than the technique I taught. My first three drafts were in my register, not his, and I now see that good feedback to a peer presenter means writing in their voice, not polishing yours. The scarce resource was register, not prose quality. I initially framed this as ``I was helping Demetro get better at presenting'' --- in fact it was ``I was learning that help which arrives in the wrong register is received as judgement, not support.'' That reframing is uncomfortable because it means my instinct to polish is not always a service.

\textbf{Feedback episode 2: Robin and the CENTERING timeout bug.}

\textit{Situation}: In wk19, while reviewing Robin's flight script for the \texttt{passive\_watch} integration, I noticed his \texttt{CENTERING} state had no timeout guard: if the target was lost mid-descent, the loop would hang indefinitely --- the drone would hover at altitude, burning battery, waiting for a detection that might never come.

\textit{Task}: Alert Robin to the specific bug without taking over his code.

\textit{Action}: I sent him a Teams message naming the exact variable (\texttt{consecutive\_lost}), the exact line, and the exact failure mode. I suggested a time-based timeout with fallback to ascend-and-retry. I did not write the fix. This was deliberate: I was giving him the diagnosis, not the prescription, because the prescription in his module should be in his style.

\textit{Result}: He fixed it the same afternoon. The behaviour change that mattered came a week later: Robin pulled me aside after a team call and said he had caught a similar timeout bug in a separate module --- a different state, same shape of bug --- because he was now ``looking for that shape of bug.'' The feedback had not just fixed one bug; it had installed a detection pattern.

\textit{Reflection}: The generalisable lesson is that when giving feedback to a high-performing peer, specificity is the scarce resource, not softness. I didn't need to wrap the message in compliments or soften the language; I needed to point at the exact line and the exact failure mode. Robin is a good enough engineer that ``here is the problem'' is all the feedback he needs. What I learned about myself: my instinct to ``finish debugging before asking'' is not professionalism. It is a cost I force teammates to pay for my own comfort with solitude. I should be asking the moment I notice I have been stuck for twenty minutes, not the moment frustration has built up far enough that the help-request feels validated to me. That pattern --- waiting until the frustration is ``enough'' --- is not a trait; it is a habit, and habits can be changed with mechanisms. The mechanism: a literal 20-minute timer. If I have been stuck for 20 minutes, I message a teammate, regardless of whether the problem feels ``big enough'' to justify the interruption.

\textbf{Feedback episode 3: The pilot and jargon removal (non-technical audience).}

\textit{Situation}: In wk22, I showed [PLACEHOLDER: pilot name] the first version of the \texttt{pi\_flight.py} browser dashboard --- the ground station he would use to monitor and command the drone during autonomous flights.

\textit{Task}: Make the dashboard usable by a pilot who holds an RC transmitter in both hands and has no time to read labels.

\textit{Action}: Two adaptations. First, the physical layout: I proposed a minimal four-button set (Y / N / M / L) mapped to the RC thumb-reach pattern, replacing a 12-button grid that required reading labels. I wired an emergency RTL into the override guard so mode 9 (LAND) from the RC switch would be honoured without dashboard interaction. Second, the lexical adaptation: my first draft used phrases like ``geofence repulsor active within 10\,m buffer'' which I rewrote as ``the drone pushes itself away from the red fence if it gets within 10\,m.'' The \texttt{FIELD\_QUICK\_REF.md} reference card used only words the pilot had used in conversation --- no acronyms, no function names, no config variables.

\textit{Result}: The rewritten card was read aloud by the pilot in one pass without a pause. During the field-day dry runs, the pilot triggered the abort button three times without looking at the screen. The moment my frame shifted from ``designing for me-as-operator'' to ``designing for a pilot who does not have to notice the UI.''

\textit{Reflection}: This episode is where M17.c (non-technical audience) lives for me. The pilot is not an engineer. He does not read Python. His operational context is completely different from mine: he holds a transmitter, looks at the sky, and needs to make split-second decisions about whether to override the autonomous system. The adaptations I made map directly to Graham's (2026) cultural competence framework: I paused, considered my own lens (``I design for engineers''), took his perspective (``he designs for physical safety''), and adapted. The jargon removal was the easy part. The harder part was restructuring the information hierarchy: my first draft put detection confidence at the top of the dashboard because that is what I care about. His hierarchy was: (1) is the drone under my control? (2) is it inside the fence? (3) has it found something? I redesigned the dashboard to match his hierarchy, not mine. In Shore et al.'s (2011) inclusion framework, this was a move from Assimilation (``use my interface'') toward Inclusion (``use an interface designed for your role'').

The A/B comparison (M17.d): the first dashboard version (12 buttons, engineering labels, detection confidence prominent) was tested on 2026-03-30. The second version (4 buttons, plain language, safety status prominent) was tested on 2026-03-31. The pilot's reaction time to the abort button decreased from approximately 3 seconds (search + read + confirm) to under 1 second (muscle memory). The method that worked was the one designed for the audience, not the one designed for the author. I would use method B in every future project involving a non-technical operator, and I would design it with the operator present, not by guessing from my desk.

\textbf{Infrastructure for team memory.} The quieter half of my support work was building infrastructure for the team's memory rather than its code:

\begin{itemize}[nosep]
\item \texttt{brain-dump.md} --- append-only capture of every idea and decision, so nothing was lost between meetings
\item \texttt{MEMORY.md} --- a persistent project state file that any team member (or any new session) could read to catch up in five minutes instead of re-reading 800 lines of CLAUDE.md
\item The blueprint system (\texttt{docs/main\_blueprint.md}, \texttt{docs/simple\_simulator\_blueprint.md}, etc.) --- line-range maps for every file over 500 lines, so you could find the function you needed without reading the whole file
\item \texttt{FIELD\_QUICK\_REF.md} --- copy-paste-ready terminal commands for flight day, written on the first cancelled field day when I realised the team would be standing in a field with no internet access
\item \texttt{docs/FLIGHT\_DAY\_CHECKLIST.md} --- a printable, single-page, follow-top-to-bottom checklist so the team could run through pre-flight safety without me standing over their shoulder
\item 127 unit tests across 6 files (\texttt{tests/unit/}) --- not required by the brief, not counted in any deliverable, but present because I believe test infrastructure is a form of team documentation
\end{itemize}

These artefacts were not glamorous. They did not train models or fly drones. But they were the substrate that let the team function when I was not in the room, and that --- making the team function without me --- is, I now believe, the highest-impact work an engineer can do on a team project.

% ============================================================
% CLOSING
% ============================================================

\subsection{What I leave believing}

I came to this project as a Ukrainian-trained engineer whose reflex under failure was ``one person delivers.'' I leave believing its near-inversion: \textit{the engineer who delivers alone is one recruitment away from being the single point of failure in whatever she builds.}

The wk17 FSM conflict and Robin's wk20 unaided handoff told me, in different registers, that the same act which delivered the report was the act that kept my teammates standing at the edge of their own project. I now see this is because, as Lencioni puts it, ``it's generally not a lack of resources or technical skill that prevents most teams from succeeding, but is often more related to a set of basic human behaviours that you yourself can support the team in strengthening.'' My resources and technical skill were not the limiting factor. My inability to share ownership was.

Applying Tuckman's model to the full arc: we Formed briefly in wk12, Stormed through wk17--20, reached a fragile Norming by wk21 when the handoff contracts held and the interfaces worked, and arguably never sustained true Performing because repeated external disruptions (weather, hardware failures) threw us back into Storming. We are now in Tuckman and Jensen's (1977) Transforming stage: the project is ending, and I feel the ``sense of loss'' they describe --- not because the project was a success, but because I can now see the team we could have been if I had invested in trust-building from wk12 instead of trust-avoidance through code.

Graham's cultural competence lecture offers one more frame. On Bennett's (1986) DMIS, I started this project at Minimisation: ``we're all engineers, the technical work is what matters.'' I have moved, imperfectly, toward Acceptance: recognising that my Ukrainian reflex of individual ownership, Robin's preference for simpler interfaces, Edward's insistence on hardware verification, and Demetro's need for spoken-register collaboration are all valid approaches that the team needs to include, not assimilate. The way I'll know I've reached Adaptation --- the next stage --- is whether I can adjust my working style to a teammate's preferences \textit{before} a conflict forces me to. On this project, every adjustment I made was reactive. The aspiration for the next one is to be proactive.

Falsifiable internship signal: in my first six weeks, does at least one teammate modify a module I own without asking permission? If yes --- even once --- I have started to become an engineer whose work enables other people's work. If no, I am still the person who writes 2508-line files that only she can read, and I will need to go back to the mechanisms and ask why.

I leave this project believing that the thing I need to develop most is not a skill but a tolerance: the tolerance to watch someone else edit my code, make it worse by my standards, and call that progress. Because a module that two people can maintain at 80\% quality is more valuable than a module that one person maintains at 100\%. That is an engineering claim, not a feelings claim, and I believe it is true. The evidence: the one file Robin could edit (\texttt{passive\_watch\_2.py}, via the HTTP contract) was the one file that worked on flight day without my intervention. Everything I protected was everything that needed me. That is not a legacy I want to repeat.

What remains unresolved: I do not yet know how to build trust quickly in a team of strangers. The mechanisms I have proposed (pair-programming budget, named-beneficiary test, documentation walkthrough) are process scaffolds, not relationship tools. The relationship part --- the part Lencioni says is the foundation of everything --- I still find uncomfortable, and I suspect I will continue to find it uncomfortable for years. But I know, now, that the discomfort is data, not a deficiency, and that ignoring it is more expensive than sitting with it. That is the lesson I will carry.
